\documentclass[11pt,a4paper,modern]{FFExercises}
\usepackage[english,greek]{babel}
\usepackage[utf8]{inputenc}
\usepackage{nimbusserif}
\usepackage[T1]{fontenc}
\usepackage{amsmath}
\let\myBbbk\Bbbk
\let\Bbbk\relax
\usepackage[amsbb,subscriptcorrection,zswash,mtpcal,mtphrb,mtpfrak]{mtpro2}
\usepackage{graphicx,multicol,multirow,enumitem,tabularx,mathimatika,gensymb,venndiagram,hhline,longtable,tkz-euclide,fontawesome5,eurosym,tcolorbox,tabularray,tikzpagenodes,relsize,siunitx,eurosym}
\definecolor{maincolor}{HTML}{aa1212} % Bright Orange accent
\definecolor{secondarycolor}{HTML}{ff9900} % Amber/Yellow accent
\usetikzlibrary{calc}
\usetikzlibrary{positioning}
\tcbuselibrary{skins,theorems,breakable}
\renewcommand{\textstigma}{\textsigma\texttau}
\renewcommand{\textdexiakeraia}{}
\usepackage{svg}
\ekthetesdeiktes
\begin{document}

\titlos{Μαθηματικά}{Γ' ΕΠΑΛ}{Μέτρα Διασποράς}

\paragraph{Ορισμοί Θεωρία}

\askhsh
Να χαρακτηρίσετε τις παρακάτω προτάσεις ως Σωστές (\textbf{Σ}) ή Λανθασμένες (\textbf{Λ}):
\begin{alist}
    \item Το εύρος $R$ ενός δείγματος ισούται με τη διαφορά της μικρότερης από τη μεγαλύτερη παρατήρηση.
    \item Η διακύμανση $s^2$ εκφράζει τον μέσο όρο των τετραγώνων των αποκλίσεων των παρατηρήσεων από τη μέση τιμή.
    \item Η τυπική απόκλιση $s$ είναι πάντα μη αρνητικός αριθμός.
    \item Αν όλες οι παρατηρήσεις ενός δείγματος είναι ίσες, τότε η τυπική απόκλιση είναι ίση με 1.
    \item Ο συντελεστής μεταβολής $\eng{CV}$ είναι αδιάστατο μέγεθος (δεν έχει μονάδες μέτρησης).
    \item Αν πολλαπλασιάσουμε όλες τις τιμές ενός δείγματος με τον αριθμό $-2$, τότε η τυπική απόκλιση πολλαπλασιάζεται επίσης με $-2$.
    \item Ένα δείγμα τιμών μιας μεταβλητής χαρακτηρίζεται ομοιογενές αν ο συντελεστής μεταβολής ξεπερνά το $10\%$.
    \item Αν προσθέσουμε σε όλες τις παρατηρήσεις ενός δείγματος τον αριθμό 5, το εύρος $R$ παραμένει αμετάβλητο.
    \item Η διακύμανση $s^2$ μετριέται στις ίδιες μονάδες με τη μεταβλητή.
    \item Όσο μικρότερη είναι η τυπική απόκλιση, τόσο μεγαλύτερη διασπορά παρουσιάζουν οι τιμές γύρω από τη μέση τιμή.
\end{alist}

\askhsh
Να χαρακτηρίσετε τις παρακάτω προτάσεις ως Σωστές (\textbf{Σ}) ή Λανθασμένες (\textbf{Λ}):
\begin{alist}
    \item Το εύρος εξαρτάται μόνο από τις ακραίες παρατηρήσεις του δείγματος.
    \item Αν η τυπική απόκλιση ενός δείγματος είναι μηδέν, τότε όλες οι παρατηρήσεις είναι ίσες μεταξύ τους.
    \item Ο συντελεστής μεταβολής $\eng{CV}$ χρησιμοποιείται για τη σύγκριση της διασποράς δύο δειγμάτων που έχουν διαφορετικές μονάδες μέτρησης.
    \item Αν η μέση τιμή ενός δείγματος είναι $\bar{x} = 0$, τότε ο συντελεστής μεταβολής $\eng{CV}$ δεν ορίζεται.
    \item Η διακύμανση ισούται με το τετράγωνο της τυπικής απόκλισης.
    \item Αν διπλασιάσουμε όλες τις τιμές ενός δείγματος (θετικές τιμές), τότε ο συντελεστής μεταβολής διπλασιάζεται.
    \item Οι τιμές της τυπικής απόκλισης $s$ και της διακύμανσης $s^2$ είναι πάντα ομόσημες.
    \item Σε μια κανονική κατανομή, το εύρος περιλαμβάνει όλες τις παρατηρήσεις.
    \item Αν $s=2$, τότε η διακύμανση είναι $s^2=4$.
    \item Αν σε ένα δείγμα η μέση τιμή ισούται με την τυπική απόκλιση, τότε $\eng{CV} = 100\%$.
\end{alist}

\paragraph{Πολλαπλής Επιλογής}

\askhsh
Να επιλέξετε τη σωστή απάντηση στις παρακάτω ερωτήσεις:
\begin{alist}
    \item Ποιο από τα παρακάτω ΔΕΝ είναι μέτρο διασποράς;
    \begin{alist}
        \item Εύρος
        \item Τυπική απόκλιση
        \item Διάμεσος
        \item Διακύμανση
    \end{alist}
    
    \item Η τυπική απόκλιση $s$ ενός δείγματος μετράει:
    \begin{alist}
        \item Τη θέση των παρατηρήσεων
        \item Τη συγκέντρωση των παρατηρήσεων γύρω από τη διάμεσο
        \item Τη διασπορά των παρατηρήσεων γύρω από τη μέση τιμή
        \item Το άθροισμα των παρατηρήσεων
    \end{alist}
    
    \item Αν η διακύμανση ενός δείγματος είναι $s^2 = 25$, τότε η τυπική απόκλιση είναι:
\begin{multicols}{4}
    \begin{alist}
        \item 5
        \item 25
        \item 625
        \item $\pm 5$
    \end{alist}
\end{multicols}
    
    \item Ο συντελεστής μεταβολής $\eng{CV}$ δίνεται από τον τύπο:
\begin{multicols}{4}
    \begin{alist}[leftmargin=4mm]
        \item $\dfrac{s}{\bar{x}^2}$
        \item $s \cdot \bar{x}$
        \item $\dfrac{\bar{x}}{s}$
        \item $\dfrac{s}{|\bar{x}|}$
    \end{alist}
\end{multicols}
    
    \item Ένα δείγμα θεωρείται ομοιογενές όταν:
    \begin{alist}
        \item $\eng{CV} > 10\%$
        \item $\eng{CV} < 10\%$
        \item $s > 10$
        \item $\bar{x} > 10$
    \end{alist}
\end{alist}

\askhsh
Να επιλέξετε τη σωστή απάντηση στις παρακάτω ερωτήσεις:
\begin{alist}
    \item Αν κάθε παρατήρηση ενός δείγματος αυξηθεί κατά 3, τότε η τυπική απόκλιση:
    \begin{alist}
        \item Αυξάνεται κατά 3
        \item Αυξάνεται κατά 9
        \item Παραμένει σταθερή
        \item Μειώνεται κατά 3
    \end{alist}
    
    \item Αν κάθε παρατήρηση ενός δείγματος πολλαπλασιαστεί με το 2, τότε το εύρος $R$:
    \begin{alist}
        \item Διπλασιάζεται
        \item Τετραγωνίζεται
        \item Παραμένει σταθερό
        \item Αυξάνεται κατά 2
    \end{alist}
    
    \item Το εύρος των τιμών $2, 5, 1, 9, 4$ είναι:
\begin{multicols}{4}
    \begin{alist}
        \item 7
        \item 8
        \item 9
        \item 5
    \end{alist}
\end{multicols}
        
    \item Σε ένα δείγμα με $\bar{x}=20$ και $s=2$, ο συντελεστής μεταβολής είναι:
\begin{multicols}{4}
    \begin{alist}
        \item $20\%$
        \item $2\%$
        \item $10\%$
        \item $5\%$
    \end{alist}
\end{multicols}
    
    \item Αν η τυπική απόκλιση ισούται με 0, τότε:
    \begin{alist}
        \item Όλες οι παρατηρήσεις είναι 0
        \item Η μέση τιμή είναι 0
        \item Όλες οι παρατηρήσεις είναι ίσες
        \item Το δείγμα είναι κενό
    \end{alist}
\end{alist}

\paragraph{Υπολογιστικές Ασκήσεις (Εύρος - Διακύμανση - Τυπική Απόκλιση)}

\askhsh
Να βρείτε το εύρος $R$ των παρακάτω δειγμάτων:
\begin{multicols}{2}
\begin{alist}
    \item $5, 12, 3, 8, 9$
    \item $10, 10, 10, 10$
    \item $-2, 5, 0, -5, 3$
    \item $1.5, 2.3, 0.8, 4.1$
\end{alist}
\end{multicols}

\askhsh
Δίνεται το δείγμα τιμών: $4, 6, 8, 10, 12$.
\begin{alist}
    \item Να υπολογίσετε τη μέση τιμή $\bar{x}$.
    \item Να υπολογίσετε το εύρος $R$.
    \item Να υπολογίσετε τη διακύμανση $s^2$.
    \item Να υπολογίσετε την τυπική απόκλιση $s$.
\end{alist}

\askhsh
Δίνεται το δείγμα τιμών: $3, 7, 3, 7$.
\begin{alist}
    \item Να υπολογίσετε τη μέση τιμή.
    \item Να υπολογίσετε την τυπική απόκλιση.
\end{alist}

\askhsh
Οι βαθμοί 5 μαθητών σε ένα διαγώνισμα είναι: \[10, 12, 14, 16, 18.\]
\begin{alist}
    \item Να βρείτε τη μέση τιμή και την τυπική απόκλιση της βαθμολογίας.
    \item Αν όλοι οι βαθμοί αυξηθούν κατά 2 μονάδες, να βρείτε τη νέα μέση τιμή και τη νέα τυπική απόκλιση.
\end{alist}

\askhsh
Οι ημερήσιες πωλήσεις (σε τεμάχια) ενός καταστήματος για μια εβδομάδα (6 ημέρες) ήταν:
$$10, 15, 10, 20, 25, 10$$
\begin{alist}
    \item Να βρείτε το εύρος των πωλήσεων.
    \item Να υπολογίσετε τη μέση τιμή.
    \item Να υπολογίσετε τη διακύμανση.
\end{alist}

\askhsh
Δίνονται οι παρατηρήσεις: $x_1=5, x_2=9, x_3=11, x_4=15$.
\begin{alist}
    \item Να βρείτε τη μέση τιμή $\bar{x}$.
    \item Να αποδείξετε ότι $\sum_{i=1}^{4} (x_i - \bar{x}) = 0$.
    \item Να υπολογίσετε την τυπική απόκλιση $s$.
\end{alist}

\askhsh
Σε ένα δείγμα 10 παρατηρήσεων γνωρίζουμε ότι $\sum x_i = 100$ και $\sum x_i^2 = 1090$.
\begin{alist}
    \item Να υπολογίσετε τη μέση τιμή $\bar{x}$.
    \item Να υπολογίσετε τη διακύμανση $s^2$.
    \item Να βρείτε την τυπική απόκλιση $s$.
\end{alist}

\askhsh
Δίνεται το δείγμα Α: $8, 8, 8, 8, 8$ και το δείγμα Β: $6, 7, 8, 9, 10$.
Χωρίς να κάνετε πράξεις, να δικαιολογήσετε ποιο δείγμα έχει μεγαλύτερη τυπική απόκλιση. Στη συνέχεια να επιβεβαιώσετε τον ισχυρισμό σας με υπολογισμούς.\\\\
\askhsh
Αν η μέση τιμή των τετραγώνων των παρατηρήσεων ενός δείγματος είναι 50 και το τετράγωνο της μέσης τιμής τους είναι 36, να βρείτε την τυπική απόκλιση.\\\\
\askhsh
Να βρείτε άγνωστο αριθμό $x$ στο δείγμα $4, 8, x$, αν γνωρίζετε ότι η μέση τιμή του δείγματος είναι 6. Στη συνέχεια να υπολογίσετε την τυπική απόκλιση.\\\\
\askhsh
Δίνεται δείγμα με 5 παρατηρήσεις και μέση τιμή $\bar{x} = 10$. Αν το άθροισμα των τετραγώνων των αποκλίσεων από τη μέση τιμή είναι $\sum (x_i - \bar{x})^2 = 80$, να βρείτε τη διακύμανση και την τυπική απόκλιση.\\\\
\askhsh
Ένας μαθητής έχει βαθμούς $12, 14, 15, 17$. Ποιον βαθμό πρέπει να πάρει στο επόμενο μάθημα ώστε η μέση τιμή του να γίνει 15; Με τον νέο βαθμό, ποια θα είναι η τυπική απόκλιση?\\\\
\askhsh
Αποδείξτε ότι αν σε ένα δείγμα προσθέσουμε τον αριθμό $c$ σε όλες τις παρατηρήσεις, η διακύμανση παραμένει αμετάβλητη. (Υπόδειξη: Χρησιμοποιήστε τον τύπο της διακύμανσης και την ιδιότητα της μέσης τιμής).\\\\
\askhsh
Υπολογίστε τη διακύμανση των αριθμών: $1, 2, 3, 4, 5$.\\\\
\askhsh
Βρείτε την τυπική απόκλιση των αριθμών: 
\[101, 102, 103, 104, 105\]

\paragraph{Πίνακες Συχνοτήτων - Απλά Δεδομένα}

\askhsh
Δίνεται ο παρακάτω πίνακας συχνοτήτων:
\begin{center}
\begin{mytblr}{
    colspec={cc},
    row{1}={maincolor, fg=white},
    cell{2-Z}{1-2}={maincolor!10}
}
$x_i$ & $\nu_i$ \\
1 & 2 \\
2 & 5 \\
3 & 2 \\
4 & 1 \\
\end{mytblr}
\end{center}
\begin{alist}
    \item Να υπολογίσετε τη μέση τιμή $\bar{x}$.
    \item Να υπολογίσετε τη διακύμανση $s^2$.
    \item Να βρείτε την τυπική απόκλιση $s$.
\end{alist}

\askhsh
Οι βαθμοί 10 μαθητών σε ένα τεστ είναι:
$$ 10, 10, 12, 12, 12, 15, 15, 18, 18, 20 $$
\begin{alist}
    \item Να κατασκευάσετε πίνακα συχνοτήτων.
    \item Να υπολογίσετε τη μέση τιμή και την τυπική απόκλιση.
\end{alist}

\askhsh
Δίνεται ο πίνακας κατανομής συχνοτήτων μιας μεταβλητής $X$:
\begin{center}
\begin{mytblr}{
    colspec={cc},
    row{1}={maincolor, fg=white},
    cell{2-Z}{1-2}={maincolor!10}
}
$x_i$ & $\nu_i$ \\
10 & 4 \\
12 & 3 \\
18 & 2 \\
20 & 1 \\
\end{mytblr}
\end{center}
Να βρείτε το εύρος $R$ και την τυπική απόκλιση $s$.

\askhsh
Σε μια έρευνα μετρήθηκε ο αριθμός παιδιών σε 20 οικογένειες:
\begin{center}
\begin{mytblr}{
    colspec={cc},
    row{1}={maincolor, fg=white},
    cell{2-Z}{1-2}={maincolor!10}
}
Αριθμός Παιδιών ($x_i$) & Αριθμός Οικογενειών ($\nu_i$) \\
0 & 2 \\
1 & 6 \\
2 & 8 \\
3 & 4 \\
\end{mytblr}
\end{center}
Να υπολογίσετε τη μέση τιμή και τη διακύμανση του αριθμού των παιδιών.

\askhsh
Για τα δεδομένα του προηγούμενου πίνακα, αν κάθε οικογένεια αποκτήσει ένα ακόμη παιδί, ποια θα είναι η νέα μέση τιμή και η νέα τυπική απόκλιση;

\askhsh
Δίνεται ο παρακάτω πίνακας με $\bar{x}=2.5$:
\begin{center}
\begin{mytblr}{
    colspec={cc},
    row{1}={maincolor, fg=white},
    cell{2-Z}{1-2}={maincolor!10}
}
$x_i$ & $\nu_i$ \\
1 & 4 \\
2 & $\nu_2$ \\
3 & 6 \\
4 & 2 \\
\end{mytblr}
\end{center}
\begin{alist}
    \item Να βρείτε τη συχνότητα $\nu_2$.
    \item Να υπολογίσετε την τυπική απόκλιση.
\end{alist}

\askhsh
Σε έναν πίνακα συχνοτήτων έχουμε $\sum_{i=1}^k \nu_i = 50$, $\sum_{i=1}^k x_i \nu_i = 500$ και $\sum_{i=1}^k x_i^2 \nu_i = 5450$.
Να υπολογίσετε τη διακύμανση $s^2$ και το συντελεστή μεταβολής $\eng{CV}$.

\askhsh
Δίνεται ο πίνακας:
\begin{center}
\begin{mytblr}{
    colspec={cccc},
    row{1}={maincolor, fg=white},
    cell{2}{1-4}={maincolor!10}
}
$x_i$ & 8 & 10 & 12 \\
$\nu_i$ & 5 & 10 & 5 \\
\end{mytblr}
\end{center}
\begin{alist}
    \item Να παρατηρήσετε τη συμμετρία της κατανομής. Ποια είναι η μέση τιμή χωρίς πράξεις;
    \item Να υπολογίσετε την τυπική απόκλιση.
\end{alist}

\paragraph{Πίνακες Συχνοτήτων - Ομαδοποιημένα Δεδομένα}

\askhsh
Δίνεται η παρακάτω κατανομή σε κλάσεις ίσου πλάτους:
\begin{center}
\begin{mytblr}{
    colspec={cc},
    row{1}={maincolor, fg=white},
    cell{2-Z}{1-2}={maincolor!10}
}
Κλάσεις & $\nu_i$ \\
$[0, 4)$ & 3 \\
$[4, 8)$ & 5 \\
$[8, 12)$ & 2 \\
\end{mytblr}
\end{center}
\begin{alist}
    \item Να βρείτε τα κέντρα των κλάσεων $x_i$.
    \item Να υπολογίσετε τη μέση τιμή $\bar{x}$.
    \item Να υπολογίσετε την τυπική απόκλιση $s$.
\end{alist}

\askhsh
Ο χρόνος (σε λεπτά) που χρειάστηκαν 20 μαθητές για να λύσουν ένα πρόβλημα είναι:
\begin{center}
\begin{mytblr}{
    colspec={cc},
    row{1}={maincolor, fg=white},
    cell{2-Z}{1-2}={maincolor!10}
}
Χρόνος (\si{min}) & Αριθμός Μαθητών \\
$[5, 15)$ & 4 \\
$[15, 25)$ & 10 \\
$[25, 35)$ & 6 \\
\end{mytblr}
\end{center}
Να υπολογίσετε τη διακύμανση του χρόνου επίλυσης.

\askhsh
Μετρήσαμε το βάρος 50 αντικειμένων και τα ομαδοποιήσαμε σε κλάσεις πλάτους $c=10$.
Αν $\bar{x} = 40$ και $\sum x_i^2 \nu_i = 85000$, να βρείτε την τυπική απόκλιση.

\askhsh
Σε μια ομαδοποιημένη κατανομή έχουμε:
$$ \sum_{i=1}^4 \nu_i = 10, \quad \sum_{i=1}^4 x_i \nu_i = 100, \quad \sum_{i=1}^4 (x_i - \bar{x})^2 \nu_i = 90 $$
Να βρείτε το συντελεστή μεταβολής $\eng{CV}$.

\askhsh
Δίνεται η κατανομή:
\begin{center}
\begin{mytblr}{
    colspec={cc},
    row{1}={maincolor, fg=white},
    cell{2-Z}{1-2}={maincolor!10}
}
Κλάσεις & $\nu_i$ \\
$[10, 20)$ & 20 \\
$[20, 30)$ & 30 \\
$[30, 40)$ & 10 \\
\end{mytblr}
\end{center}
Να υπολογίσετε την τυπική απόκλιση.

\askhsh
Για την κατανομή της προηγούμενης άσκησης, αν όλα τα δεδομένα αυξηθούν κατά $5\%$, ποια θα είναι η νέα τυπική απόκλιση και το νέο $\eng{CV}$;

\askhsh
Δίνονται δύο κατανομές συχνοτήτων Α και Β με την ίδια μέση τιμή.
Στην κατανομή Α οι συχνότητες είναι συγκεντρωμένες κοντά στη μέση τιμή.
Στην κατανομή Β οι συχνότητες είναι απλωμένες στις ακραίες κλάσεις.
Ποια κατανομή έχει μεγαλύτερη τυπική απόκλιση; Δικαιολογήστε.

\paragraph{Συντελεστής Μεταβολής (\eng{CV}) - Ομοιογένεια}

\askhsh
Σε ένα δείγμα με $\bar{x}=40$ και $s=2$:
\begin{alist}
    \item Να υπολογίσετε τον συντελεστή μεταβολής $\eng{CV}$.
    \item Να εξετάσετε αν το δείγμα είναι ομοιογενές.
\end{alist}

\askhsh
Δίνεται δείγμα με $\nu=20$, $\bar{x}=15$ και διακύμανση $s^2=9$.
\begin{alist}
    \item Να βρείτε την τυπική απόκλιση.
    \item Να υπολογίσετε το $\eng{CV}$. Είναι το δείγμα ομοιογενές?
\end{alist}
\askhsh
Οι μισθοί σε μια εταιρεία έχουν μέση τιμή $1200$\euro\ και τυπική απόκλιση $150$\euro. Είναι τα μισθολογικά δεδομένα ομοιογενή?\\\\
\askhsh
Σε ένα διαγώνισμα Μαθηματικών, η μέση βαθμολογία ήταν 14 και η τυπική απόκλιση 2.
Σε ένα διαγώνισμα Φυσικής, η μέση βαθμολογία ήταν 16 και η τυπική απόκλιση 3.
Ποιο από τα δύο δείγματα βαθμολογιών παρουσιάζει μεγαλύτερη σχετική διασπορά (ανομοιογένεια)?\\\\
\askhsh
Ένα δείγμα έχει $\eng{CV}=20\%$ και τυπική απόκλιση $s=4$. Να βρείτε τη μέση τιμή του.\\\\
\askhsh
Σε ένα δείγμα θετικών παρατηρήσεων, το $\eng{CV}$ είναι $50\%$. Αν διπλασιάσουμε κάθε παρατήρηση, ποιο θα είναι το νέο $\eng{CV}$; Δικαιολογήστε την απάντησή σας.\\\\
\askhsh
Εξετάστε την ομοιογένεια των παρακάτω δειγμάτων:
\begin{alist}
    \item Δείγμα Α: $\bar{x}=50, s=4$.
    \item Δείγμα Β: $\bar{x}=80, s=10$.
\end{alist}

\askhsh
Ένα δείγμα έχει μέση τιμή 10 και τυπική απόκλιση 0,8. Αν προσθέσουμε σε κάθε παρατήρηση τον αριθμό $c=10$, θα είναι το νέο δείγμα ομοιογενές;

\askhsh
Δίνεται το δείγμα $5, 9, 11, 15$.
\begin{alist}
    \item Να υπολογίσετε τη μέση τιμή και την τυπική απόκλιση.
    \item Να εξετάσετε αν το δείγμα είναι ομοιογενές.
\end{alist}

\askhsh
Για ποιες τιμές της μέσης τιμής $\bar{x}$ (με δεδομένο $s=3$) ένα δείγμα είναι ομοιογενές;

\askhsh
Ένα δείγμα Α έχει $\bar{x}_A = 20, s_A = 2$. Ένα δείγμα Β έχει $\bar{x}_B = 40, s_B = 6$.
\begin{alist}
    \item Ποιο δείγμα έχει μεγαλύτερη απόλυτη διασπορά;
    \item Ποιο δείγμα έχει μεγαλύτερη σχετική διασπορά;
\end{alist}

\askhsh
Σε μια ομάδα μπάσκετ το μέσο ύψος είναι $200$ \si{cm} με τυπική απόκλιση $10$ \si{cm}. Σε μια ομάδα ποδοσφαίρου το μέσο ύψος είναι $180$ \si{cm} με τυπική απόκλιση $12$ \si{cm}. Ποια ομάδα παρουσιάζει μεγαλύτερη ομοιογένεια ως προς το ύψος?\\\\
\askhsh
Δίνεται δείγμα 4 θετικών αριθμών με μέση τιμή 10. Αν το \eng{CV} είναι $20\%$, να βρείτε τη διακύμανση του δείγματος.\\\\
\askhsh
Αν σε ένα δείγμα ισχύει $s = \frac{1}{2} \bar{x}$, πόσο είναι το \eng{CV}; Είναι το δείγμα ομοιογενές?\\\\
\askhsh
Το βάρος 100 αυγών έχει μέση τιμή 60\si{g} και τυπική απόκλιση 6\si{g}.
\begin{alist}
    \item Είναι το δείγμα ομοιογενές;
    \item Αν αφαιρέσουμε το κέλυφος (βάρους 5\si{g} από κάθε αυγό), ποια θα είναι η νέα μέση τιμή και η νέα τυπική απόκλιση;
    \item Θα αλλάξει η ομοιογένεια του δείγματος;
\end{alist}

\askhsh
Δίνονται δύο τμήματα Γ1 και Γ2.
Στο Γ1: $\bar{x}_1 = 15, s_1 = 2$.
Στο Γ2: $\bar{x}_2 = 14, s_2 = 3$.
Ποιο τμήμα έχει πιο ομοιογενή βαθμολογία?\\\\
\askhsh
Σε έναν πίνακα συχνοτήτων, η μέση τιμή είναι 20 και η τυπική απόκλιση 4.
Αν πολλαπλασιάσουμε όλες τις συχνότητες $\nu_i$ επί 2 (δηλαδή έχουμε το ίδιο δείγμα 2 φορές), αλλάζει η ομοιογένεια?\\\\
\paragraph{Συνδυαστικά Θέματα - Κριτική Σκέψη}

\askhsh
Δίνεται η συνάρτηση $f(x) = x^2 - 4x + 6$.
\begin{alist}
    \item Να βρείτε το ελάχιστο της συνάρτησης.
    \item Θεωρούμε δείγμα με 3 παρατηρήσεις οι οποίες είναι οι ρίζες της εξίσωσης $(x-2)(x^2 - 5x + 6) = 0$.
    \begin{alist}
        \item Να βρείτε τις παρατηρήσεις.
        \item Να υπολογίσετε τη μέση τιμή και την τυπική απόκλιση.
        \item Να εξετάσετε την ομοιογένεια.
    \end{alist}
\end{alist}

\askhsh
Σε ένα εργοστάσιο, δύο μηχανές Α και Β παράγουν εξαρτήματα.
Η μηχανή Α παράγει εξαρτήματα με μέσο μήκος 10\si{cm} και $s=0.1$\si{cm}.
Η μηχανή Β παράγει εξαρτήματα με μέσο μήκος 5\si{cm} και $s=0.08$\si{cm}.
Ποια μηχανή είναι πιο «σταθερή» (ακριβής) στην παραγωγή της?\\\\
\askhsh
Ένα δείγμα αποτελείται από 10 άντρες και 10 γυναίκες.
Το μέσο ύψος των ανδρών είναι 180\si{cm} και των γυναικών 170\si{cm}.
Μπορούμε να υπολογίσουμε την τυπική απόκλιση του συνολικού δείγματος μόνο από τους μέσους όρους; (Απάντηση: Όχι - Ερώτηση κρίσεως).
Τι επιπλέον στοιχεία θα χρειαζόμασταν?\\\\
\askhsh
\textbf{Βρείτε το λάθος:}
Ένας μαθητής υπολόγισε την τυπική απόκλιση ενός δείγματος και βρήκε $s = -2$.
Πού βρίσκεται το σφάλμα?
\begin{alist}
    \item Στις πράξεις?
    \item Στον τύπο?
    \item Στην κατανόηση της έννοιας?
\end{alist}
Δικαιολογήστε.
\askhsh
Δίνονται τρία δείγματα με την ίδια μέση τιμή $\bar{x} = 50$.
Το δείγμα Α έχει εύρος $R=10$.
Το δείγμα Β έχει εύρος $R=20$.
Το δείγμα Γ έχει εύρος $R=5$.
Ποιο δείγμα περιμένετε να έχει τη μικρότερη τυπική απόκλιση? (Υποθέτοντας κανονική κατανομή).

\askhsh
Δίνεται δείγμα 5 αριθμών με μέση τιμή 4 και $s^2=0$.
Ποιοι είναι οι αριθμοί;

\askhsh
Δίνεται το δείγμα $x, x+2, x+4, x+6$.
Να δείξετε ότι η διακύμανση είναι ανεξάρτητη του $x$.

\askhsh
Σε μια κανονική κατανομή, το $68\%$ των παρατηρήσεων βρίσκεται στο διάστημα $(\bar{x}-s, \bar{x}+s)$.
Αν $\bar{x}=10$ και $s=2$, σε ποιο διάστημα βρίσκεται το $68\%$ των τιμών;

\askhsh
Αν σε ένα δείγμα η μέση τιμή τριπλασιαστεί (με πρόσθεση σταθεράς), πώς πρέπει να μεταβληθεί η τυπική απόκλιση ώστε το \eng{CV} να παραμείνει σταθερό;

\askhsh
Για 10 ζεύγη παρατηρήσεων $(x_i, y_i)$ ισχύει $y_i = 2x_i + 3$.
Αν $s_x = 2$, να βρείτε την $s_y$.

\end{document}
